\section{Protocol Robustness and Integrity}
\label{sec:robust}

\begin{guidance}[Protocol robustness and integrity (10 marks)]
Show \textbf{framing}, \textbf{specific NACKs for malformed input}, and
\textbf{graceful recovery} from partial or garbage data. A table that maps each
kind of bad input to your mechanism, its response and the test that checks it
works well. Also explain why the server can never block forever (bounded
loops, timeouts, no blocking waits in the main loop). Give measured numbers
where you have them, e.g. the timeout you observed.
\end{guidance}

The server refuses every bad frame with a specific NACK, discards it up to the
next terminator or the 500~ms timeout, and answers the next valid command
normally. Table~\ref{tab:threats} maps each kind of bad input to the mechanism
and the test that exercises it.

\begin{table}[t]
\centering
\caption{Bad input, the server's handling, and the test that checks it.}
\label{tab:threats}
\scriptsize
\begin{tabular}{@{}L{1.55cm}L{3.0cm}L{3.45cm}@{}}
\toprule
Bad input & Handling and response & Test \\
\midrule
Partial frame & dropped after 500~ms of silence; \texttt{NACK:TIMEOUT} & \texttt{test\_partial\_frame\_times\_out\_and\_is\_discarded}; \texttt{test\_frame\_split\_across\_writes} (a split under 500~ms is reassembled) \\
Oversized frame & discarded to terminator; one \texttt{NACK:TOOLONG} & \texttt{test\_oversized\_frame\_rejected\_and\_tail\_not\_executed} \\
Garbage bytes & non-printable byte flags the frame; \texttt{NACK:BADCMD} & \texttt{test\_garbage\_bytes\_rejected\_then\_recovers} \\
Wrong case or unknown command & exact match only; \texttt{NACK:BADCMD} & \texttt{test\_unknown\_and\_wrong\_case\_commands}; \texttt{test\_malformed\_move\_does\_not\_touch\_state}; \texttt{test\_malformed\_suffix\_is\_not\_a\_checksum} \\
Corrupted bytes & CRC-8 suffix checked before parsing; \texttt{NACK:CHECKSUM}, nothing runs (frames that carry a suffix) & \texttt{test\_wrong\_checksum\_rejected\_and\_nothing\_executed}; \texttt{test\_corrupted\_payload\_is\_detected} \\
UART errors & flags cleared in the interrupt; frame rejected with \texttt{NACK:BADCMD} & no dedicated test; first seen through the oversize test \\
\bottomrule
\end{tabular}
\end{table}

\subsection{A 500 ms timeout is far longer than any frame}
A 64-byte frame takes about 5.6~ms on the wire (10 bits per byte at 115200
baud), so a healthy sender never reaches 500~ms between bytes, and a stalled
sender is cleared within about half a second. Our test checks only that the
\texttt{NACK:TIMEOUT} arrives inside the 2~s read window of the harness; we did
not measure the exact delay.

\subsection{The checksum detects corruption the framer cannot see}
Framing only checks the shape of a frame. A flipped bit that turns \texttt{MOVE:4}
into \texttt{MOVE:5} still has the right shape, so the framer accepts it and the
server would play the wrong cell. The CRC-8 suffix closes that gap for frames that
carry it: \texttt{test\_corrupted\_payload\_is\_detected} sends exactly that
substitution with the old checksum and checks that the board does not change. A CRC
with a degree-8 polynomial detects every single-bit error and every burst of up to
8 bits. The protection is optional: a plain frame has none, and damage to the
\texttt{*} or the two digits makes a checksummed frame look plain, when it is
usually (not provably) rejected as an unknown command.

\subsection{The next terminator resynchronises the stream}
A bad frame sets \texttt{frame\_err}, the framer drops the remaining bytes,
and the first \texttt{\textbackslash n} or \texttt{\textbackslash r} clears the flag
after one NACK. The tests send a bad frame and then a valid command, and
check that the valid command gets its normal reply and that the bad frame's
tail did not run (Listing~\ref{lst:framer}).

\subsection{No path in the main loop waits without a bound}
The framer loop ends when the ring is empty, which is at most 127 bytes. Transmit
is polled with a 100~ms timeout per message. The main loop contains no
\texttt{HAL\_Delay}: the button is debounced by comparing timestamps, and the only
delays are the three LED flashes (600~ms) at start-up, before any request is
accepted. We did not enable the independent watchdog, and
\texttt{Error\_Handler} halts only if the HAL initialisation itself fails.
